\documentclass[11pt]{article}
\usepackage[english]{babel}
\usepackage[a4paper,margin=0.95in]{geometry}
\usepackage[T1]{fontenc}
\usepackage[utf8]{inputenc}
\usepackage{lmodern}
\usepackage{amsmath,amssymb,mathtools,mathrsfs}
\usepackage{microtype}
\usepackage{fancyhdr}
\setlength{\parindent}{1.5em}
\setlength{\parskip}{0.25em}
\emergencystretch=6em
\setlength{\headheight}{14pt}
\pagestyle{fancy}
\fancyhf{}
\cfoot{\thepage}
\lhead{Weber Vol. I}
\rhead{Introduction: measurable sets and ratios}
\begin{document}
\begin{center}{\Large Introduction: measurable sets and ratios}\end{center}

This very abstract way of looking at the matter gives us the assurance that the assumption of a continuous set contains no contradiction, and that such sets at least exist in the realm of thought. Geometry, and also analysis, which has always willingly attached itself to geometrical intuition, long tacitly assumed the existence of continuous sets, for example in the points of a straight line or of any other connected line-segment, as a kind of axiom. The difference between a dense and a continuous set, which lies at the basis of the distinction between commensurable and incommensurable segments, was also not missed by the ancients.\footnote{Euclid, Elements, Book X.}

We too shall in what follows not renounce the aid of geometrical intuition, and shall, for example, without hesitation regard the points of a straight line as a continuous set.

An ordered set $M$ is called measurable under the following assumptions. Addition and multiplication are generally executable in $M$, and likewise the subtraction of a smaller element from a greater one; that is, from any two elements $a,b$ (which may also be identical) a new element $a+b$ of $M$ can be derived according to a definite rule, in such a way that $a+b$ is greater than $a$ and greater than $b$, and the familiar laws of addition expressed in the formulae
\[
 a+b=b+a,
 \qquad (a+b)+c=a+(b+c)
\]
hold. Moreover, for two elements $a,c$, of which the second is the greater, a third element $b$ can be found such that $a+b=c$, which is also expressed by the sign of subtraction $b=c-a$. Thus two unequal elements of $M$ always have a definite difference. From these assumptions it follows that a sum becomes greater when one of its summands becomes greater. The repeated, say $m$-fold, addition of the same element $a$ is called multiplication, and its result is denoted by $ma$.

Finally there is added one more assumption, namely that, for every given $a$, a sufficiently high multiple $ma$ is greater than any other element $b$ given in advance. Among the elements of a measurable set there is therefore no greatest element.

In a dense measurable set there is also no least element. For if $a$ were the least and $b$ an arbitrary element, then between $b$ and $b+a$ no elements could lie; for if $b<c<b+a$, then it would follow from the definition of measurability that $a'=c-b$ would be smaller than $a$. Conversely it also follows that a measurable set in which no least element occurs is dense. For if $a$ and $a+b$ are any two elements, one need only add to $a$ an element smaller than $b$ in order to obtain an element between $a$ and $a+b$.

If a set is continuous, then the assumptions for measurability can be simplified still further, because then subtraction is a consequence of addition. For if $a$ and $c$ are two elements of a continuous ordered set $M$ in which addition exists, and if $c>a$, then one obtains a cut $(A,B)$ in $M$ by assigning to $A$ all elements $x$ for which $a+x<c$, and to $B$ those for which $a+x>c$. This cut is produced by an element $b$, for which $a+b=c$ must hold.

The natural numbers form, according to our definition, a measurable set in which a least element, namely $1$, occurs. The manifold of rational fractions is likewise measurable when addition and subtraction are explained according to the known rules of fractional calculation. Especially important, and as it were typical for the measurable sets, is the manifold of rectilinear segments, or lengths of a line, which are added simply by laying them end to end. Quantities of material, compared by means of the balance, and intervals of time, measured by the clock, also furnish examples of measurable sets. The manner of measuring does not lie in the nature of the manifold itself, but is put into it by the thinking observer; thus, for example, it would be just as admissible to understand by the sum $a+b$ of two segments $a$ and $b$ the hypotenuse of a right triangle with legs $a,b$, instead of, as is usually assumed, the segment composed from $a$ and $b$ by laying them end to end.

In order to make the continuous set $\mathfrak S$ of cuts in the manifold $\mathfrak R$ of rational fractions into a measurable one, observe first the following. If $\alpha=(A,B)$ is an element in $\mathfrak S$ and $\mu$ is an arbitrarily given rational fraction, then one can always choose an element $a'$ in $A$ so that $a'+\mu=b'$ is contained in $B$. For if one chooses two arbitrary elements $a,b$, then one can determine the natural number $m$ so that $m\mu>b-a$, hence $a+m\mu$ is contained in $B$. If then $h$ is the smallest integer for which $a+h\mu$ is contained in $B$, then $a+(h-1)\mu=a'$ is in $A$ and hence $a'+\mu=b'$ is in $B$.

We now understand by the sum $(A,B)+(A',B')=(A'',B'')$, or $\alpha+\alpha'=\alpha''$, the cut in $R$ which is obtained by assigning a rational fraction $a''$ to $A''$ only when an $a$ in $A$ and an $a'$ in $A'$ exist such that $a''\le a+a'$. In fact $(A'',B'')$ is a cut; for if $a''$ is contained in $A''$, the same holds for every smaller fraction, and there are fractions which are in $A''$ and fractions which are not in $A''$, namely fractions of the forms $a+a'$ and $b+b'$. Further, $\alpha''$ is greater than $\alpha$ and than $\alpha'$. For $A''$ first contains all of $A$, and if $a'$ is an arbitrary fraction in $A'$, one can choose $a$ in $A$ so that the element $a+a'$ of $A''$ is contained in $B$; hence $A''$ is more comprehensive than $A$. The cuts produced by rational fractions $\mu,\mu'$ give, by addition, the cut produced by $\mu+\mu'$.

We now pass to the definition of ratios, which since antiquity have been regarded as the foundation of the theory of numbers, and here follow Euclid first of all.\footnote{Elements, Book V.}

If one combines the elements of a measurable set $M$ into pairs, and regards these pairs themselves as elements, then a new manifold arises. We denote such a pair by $a:b$, or also by $a/b$; but, when $a$ and $b$ are different elements, we distinguish $a:b$ from $b:a$, and call $a$ the numerator and $b$ the denominator of $a:b$. We shall call these pairs ratios, and shall now order and make measurable this new set.

Assume first that two whole numbers $m,n$ exist such that $na=mb$, as is always the case, for example, when $M$ is the system of natural numbers, or when $a,b$ are two commensurable segments. Then, if $p,q$ are two other whole numbers, $qa=pb$ holds if and only if $mq=np$. The pair of numbers $p,q$ is completely determined by this requirement if the condition is added that $p,q$ are to be relatively prime. Then, if $h$ is an arbitrary whole number, $m=hp$, $n=hq$. In this case we call the ratio $a:b$ rational and set it equal to the rational fraction $m:n$ or $p:q$. Accordingly these rational fractions may be regarded as ratios of whole numbers.

All rational ratios equal to one another form a rational number, and the rational numbers form, like the rational fractions, an ordered, dense, and measurable manifold. The natural numbers themselves are inserted into the manifold of rational numbers if by the natural number $m$ one understands the ratio $m:1$.

We now return to any measurable manifold $M$ and take from it any two elements $a$ and $b$. If one chooses, which is always possible, two natural numbers $m,n$ so that $na>mb$, then the ratio $a:b$ is called greater than the rational ratio $m:n$, or
\[
 {a\over b}>{m\over n};
\]
and if $m:n=p:q$, then also $a:b>p:q$. Similarly, if $n'a<m'b$, then
\[
 {a\over b}<{m'\over n'}.
\]
If $a:b>m:n$, then one can insert one, and consequently also arbitrarily many, rational numbers $m_1:n_1$ such that
\[
 {a\over b}>{m_1\over n_1}>{m\over n}.
\]
For this choose an arbitrary whole number $k$ and determine the whole number $h$ so that $h(na-mb)>kb$, which is always possible. Then
\[
 {a\over b}>{hm+k\over hn}>{m\over n}.
\]
And in exactly the same way, if $n'a<m'b$ and $h'(m'b-n'a)>k'a$, one obtains
\[
 {a\over b}<{h'm'
\over h'n'+k'}<{m'
\over n'}.
\]

If now $a:b$ and $\alpha:\beta$ are any two ratios, which for brevity we shall also denote by $\epsilon,\epsilon'$, and whose elements belong to the same or to different manifolds, then two cases are possible: either no rational ratio $\mu$ lies between $\epsilon$ and $\epsilon'$, or a rational ratio lies between them.

In the first case the two ratios $\epsilon$ and $\epsilon'$ are called equal. One sees that two ratios which are equal to a third are also equal to one another; for if $\epsilon<\mu<\epsilon'$, then every other ratio $\epsilon''$ is either smaller than, equal to, or greater than $\mu$. If $\epsilon''$ is equal to $\mu$, then rational ratios lie both between $\epsilon$ and $\epsilon''$ and between $\epsilon''$ and $\epsilon'$. If, however, $\epsilon''$ is smaller than $\mu$, then $\mu$ lies between $\epsilon''$ and $\epsilon'$, and if $\epsilon''$ is greater than $\mu$, then $\mu$ lies between $\epsilon$ and $\epsilon''$; hence $\epsilon''$ cannot at the same time be equal to $\epsilon$ and equal to $\epsilon'$.

In the second case the ratios $\epsilon,\epsilon'$ are called unequal. Thus either $\epsilon<\mu<\epsilon'$ or $\epsilon>\mu'>\epsilon'$ may hold, and these two cases exclude one another, because from $\epsilon<\mu<\epsilon'$ and $\epsilon<\mu'$ it follows that $\mu<\mu'$, and consequently $\epsilon<\mu'$ results.

The relation of magnitude $2\alpha)$ or $2\beta)$ also remains unchanged when $\epsilon$ or $\epsilon'$ is replaced by an element equal to it. For if $\mu<\epsilon$ and $\mu\gtrless\epsilon'$, then a rational ratio lies between $\epsilon$ and $\epsilon'$, and $\epsilon,\epsilon'$ are not equal.

In the case $2\alpha)$ the ratio $\epsilon$ is called smaller than $\epsilon'$; in the case $2\beta)$ it is called greater than $\epsilon'$. Between two unequal ratios one can insert an arbitrary number of rational ratios.

If now we collect all ratios equal to one another, we obtain a generic concept which we designate as a number in the general sense of the word. A number is therefore a name or sign for a certain manifold whose elements are precisely the ratios equal to one of them.\footnote{The natural numbers too can be traced back in a simple and consistent way to the generic concept.} Under this concept of number are included the rational ratios, and hence also the natural numbers as the ratios $m:1$; they form the rational numbers. Numbers which do not arise from rational ratios are called irrational numbers. According to what has been developed so far, the numbers form an ordered set, and their order can be fixed if, for each number, any one of the ratios contained under it is chosen as representative.

Of two ratios with the same denominator and unequal numerators, that one is greater whose numerator is greater; and of two ratios with the same numerator and unequal denominators, that one is smaller whose denominator is greater. For if $a,a',b$ are arbitrary elements of a measurable set and $a'>a$, choose first a whole number $n$ such that $na>b$ and $n(a'-a)>b$. Then take $m$ to be the least whole number satisfying $mb>na$; then $na<mb$, but $mb<na'$. For if $mb\ge na'=na+n(a'-a)$, then $mb>na+b$; hence, contrary to the supposition, already $(m-1)b>na$. Thus
\[
 {a\over b}<{m\over n}<{a'
\over b},
\]
and hence $a:b<a':b$. In exactly the same way one proves that if $b'>b$, then $a:b>a:b'$. This theorem is also expressed by saying that a ratio increases with the numerator and decreases with increasing denominator.

From this one easily obtains the following theorem. If $a,b,c,d$ are elements of the same measurable set and $a:b=c:d$, then also $a:c=b:d$. For suppose that $a:c<b:d$; then there would have to be two whole numbers $m,n$ such that
\[
 na<mc,
 \qquad nb>md.
\]
Then, however, by the theorem just proved, $na:nb<mc:md$, and hence $a:b<c:d$, contrary to the supposition.

To this is now attached the following principal theorem. If of four quantities $a,b,c,d$ any three are given arbitrarily from a continuous measurable set, then the fourth can be determined in the same set so that $a:b=c:d$.

The theorem is an immediate consequence of the assumed continuity. For if, for example, an element $x$ of $M$ is put into $A$ or into $B$ according as $x:b$ is smaller or greater than $c:d$, then a cut is obtained, which is produced by an element $a$ satisfying $a:b=c:d$. But this theorem also holds in certain non-continuous sets, for example for the rational fractions.

It follows that as representatives of two numbers one can always choose two ratios whose elements belong to the same manifold and have the same arbitrarily chosen denominator. The addition is then defined by
\[
 {a\over c}+{b\over c}={a+b\over c}.
\]
This rule contains as a special case the addition of rational fractions; and to justify it generally it remains only to show that if $a:c=a':c'$ and $b:c=b':c'$, then also $(a+b):c=(a'+b'):c'$. We prove this by showing that if $a:c=a':c'$ and $(a+b):c>(a'+b'):c'$, then also $b:c>b':c'$ must hold. Let therefore, with $m,n$ two whole numbers,
\[
 {a+b\over c}>{m\over n}>{a'+b'
\over c'}.
\]
Then $n(a+b)>mc$ and hence
\[
 {b\over c}>{mc-na\over nc}.
\]
On the other hand, from the supposition $a:c=a':c'$ it follows easily that
\[
 {mc-na\over nc}={mc'-na'\over nc'},
\]
and therefore $b:c>b':c'$, as was required to be proved.

It is thus shown that the numbers, as we have defined them, also form a measurable set. If $a$ and $c$ are taken from a continuous manifold, then, with fixed $c$, the ratios $a:c$ also form a continuous set; hence, since continuous sets exist at all, it follows that the numbers too form a continuous set.
\end{document}
